% ============================================
% LUC-02 (EN)
% English translation — SPEC-935-R355
% ============================================
\section*{\textbf{LUC-02} --- The Archive}
\markboth{LUC-02 --- The Archive}{}

\noindent \textit{The letter from 1979 to 2026}

\noindent \textbf{Public Archive of Minas Gerais --- Belo Horizonte}
\textbf{February 2026}

\vspace{0.5em}\noindent\rule{\textwidth}{0.5pt}\vspace{0.5em}

\noindent The Public Archive of Minas Gerais stands in a gray building in the center of Belo Horizonte. Lúcia knew the place --- she had spent hours there for other research. But she had never seen the back room.

\noindent The director, Dr. Ronaldo Torres, was waiting for her at the door. He was a thin man, with thick glasses and restless hands. He greeted her with a quick handshake.

\noindent \textit{\textquotedblleft{}Dr. Lúcia. Thank you for coming so quickly.\textquotedblright{}}

\noindent \textit{\textquotedblleft{}You said it was urgent.\textquotedblright{}}

\noindent \textit{\textquotedblleft{}It is.\textquotedblright{}\textit{ He looked around. }\textquotedblleft{}Let us go to my office.\textquotedblright{}}

\noindent The director's office was small, full of folders and books. On the desk, a package wrapped in black cloth and, beside it, a rusty key on a sheet of tissue paper.

\noindent \textit{\textquotedblleft{}It arrived last week\textquotedblright{}\textit{, said Ronaldo. }\textquotedblleft{}A woman. She left no name. She said only: 'This belongs to the archive. Dr. Heitor Oliveira asked for it to be returned.'\textquotedblright{}}

\noindent \textit{\textquotedblleft{}Heitor Oliveira has been missing since 1979.\textquotedblright{}}

\noindent \textit{\textquotedblleft{}I know.\textquotedblright{}}

\noindent \textit{\textquotedblleft{}How could he ask for a diary to be returned almost fifty years later?\textquotedblright{}}

\noindent Ronaldo did not answer. He took off his glasses and cleaned the lenses.

\noindent \textit{\textquotedblleft{}Do you want to see the diary?\textquotedblright{}}

\noindent \textit{\textquotedblleft{}Yes.\textquotedblright{}}

\noindent \textit{\textquotedblleft{}Before that, I need to show you something.\textquotedblright{}}

\noindent He opened a folder on the desk. Inside, there were photographs. Old. Yellowed. The first showed a thin boy, with skinny arms and a yellow right eye. He stood in front of a brick wall, with a number hanging around his neck: \textit{1,260}.

\begin{figure}[H]
  \centering
  \includegraphics[width=0.55\textwidth]{foto_01_paciente1260_sepia.png}
  \caption*{\footnotesize\textbf{Photograph 1} --- \textit{Joaquim Antônio Correia --- Patient 1,260. 1917.}}
\end{figure}

\noindent \textit{\textquotedblleft{}Joaquim\textquotedblright{}}, said Lúcia.

\noindent \textit{\textquotedblleft{}Yes. This is the only photo of him. It was taken upon his arrival at the Colony, in 1917.\textquotedblright{}}

\noindent \textit{\textquotedblleft{}And the others?\textquotedblright{}}

\noindent Ronaldo turned the photos over. The second showed a man --- a young physician, in a lab coat, with deep dark circles and a forced smile.

\begin{figure}[H]
  \centering
  \includegraphics[width=0.55\textwidth]{foto_02_dr_heitor_sepia.png}
  \caption*{\footnotesize\textbf{Photograph 2} --- \textit{Dr. Heitor Oliveira. Colony archive, 1979.}}
\end{figure}

\noindent \textit{\textquotedblleft{}Dr. Heitor Oliveira. Badge photo, 1979.\textquotedblright{}}

\noindent The third photo made Lúcia stop.

\noindent It was a photo of an older man. Much older. White hair, wrinkled skin, dull eyes. He sat on a park bench, with a book on his lap.

\begin{figure}[H]
  \centering
  \includegraphics[width=0.55\textwidth]{foto_03_homem_livro_sepia.png}
  \caption*{\footnotesize\textbf{Photograph 3} --- \textit{Unidentified man. Found inside the diary.}}
\end{figure}

\noindent \textit{\textquotedblleft{}Who is he?\textquotedblright{}} asked Lúcia.

\noindent \textit{\textquotedblleft{}I do not know\textquotedblright{}\textit{, said Ronaldo. }\textquotedblleft{}This photo was inside the diary. Together with a letter.\textquotedblright{}}

\noindent \textit{\textquotedblleft{}Letter?\textquotedblright{}}

\noindent \textit{\textquotedblleft{}Addressed to you.\textquotedblright{}}

\noindent Ronaldo picked up a brown envelope. It was sealed. Written by hand:

\noindent \textit{\textquotedblleft{}For Dr. Lúcia Mendes --- Personal\textquotedblright{}}

\noindent Lúcia pointed at the rusty key on the desk. \textit{\textquotedblleft{}And the key?\textquotedblright{}}

\noindent Ronaldo picked up the key carefully. \textit{\textquotedblleft{}It came with it. I do not know if it is from the diary, the archive, or some Colony cell. We have tried it on every door in the building. It fits none.\textquotedblright{}}

\noindent \textit{\textquotedblleft{}Why was it sent along?\textquotedblright{}}

\noindent \textit{\textquotedblleft{}The woman who delivered the package said: 'Dr. Oliveira asked me to tell you that the key does not open doors. It opens the story.'\textquotedblright{}}

\noindent Lúcia touched the key. The metal was cold --- colder than it should be in a closed room.

\noindent \textit{\textquotedblleft{}Keep it with the diary\textquotedblright{}\textit{, she said. }\textquotedblleft{}It may be important.\textquotedblright{}}

\noindent \textit{\textquotedblleft{}How did they know I was going to request the archive?\textquotedblright{}}

\noindent \textit{\textquotedblleft{}I do not know.\textquotedblright{}}

\noindent Lúcia opened the envelope. Inside, a yellowed paper, in trembling handwriting:

\vspace{0.5em}\noindent\rule{\textwidth}{0.5pt}\vspace{0.5em}

\noindent \textit{Dr. Lúcia,}

\noindent \textit{If you are reading this letter, the cycle has reached you.}

\noindent \textit{My name is Heitor Oliveira. I am writing this letter in 2026. Or rather, I am writing it in 1979 and asking for someone to deliver it in 2026.}

\noindent \textit{Do not ask how. The diary teaches.}

\noindent \textit{I read patient 1,260's diary in March 1979. Since then, I am no longer the same.}

\noindent \textit{I spent 47 years trying to understand what happened to me. Trying to understand what happened to Joaquim.}

\noindent \textit{I did not understand.}

\noindent \textit{I understood only one thing: the diary is not a record. It is a vehicle. The creature Joaquim found in the trench does not live in bodies. It lives in stories. Every time someone reads, it passes to the reader.}

\noindent \textit{Have you read this diary?}

\noindent \textit{No. Not yet.}

\noindent \textit{But you will read it.}

\noindent \textit{And when you read it, you will understand.}

\noindent \textit{Heitor Oliveira}
\textit{--- 1979 (sent to 2026)}

\vspace{0.5em}\noindent\rule{\textwidth}{0.5pt}\vspace{0.5em}

\noindent Lúcia set the letter on the desk. Her hand was trembling.

\noindent \textit{\textquotedblleft{}This is a joke\textquotedblright{}\textit{, she said. }\textquotedblleft{}Someone is manipulating the archive.\textquotedblright{}}

\noindent \textit{\textquotedblleft{}Perhaps\textquotedblright{}\textit{, said Ronaldo. }\textquotedblleft{}But there is one more thing.\textquotedblright{}}

\noindent \textit{\textquotedblleft{}What?\textquotedblright{}}

\noindent \textit{The director pointed at the black cloth package. The diary.}

\noindent \textit{\textquotedblleft{}The woman who delivered it said another phrase. She said: 'Tell Dr. Lúcia that patient 1,260 sends his regards.'\textquotedblright{}}

\noindent Lúcia swallowed hard.

\noindent Joaquim Antônio Correia died in 1979.

\noindent That was a fact.

\noindent But the diary was there. Warm. Waiting.

\noindent \textit{\textquotedblleft{}I will take it.\textquotedblright{}}

\noindent \textit{\textquotedblleft{}Are you sure?\textquotedblright{}}, asked Ronaldo.

\noindent \textit{\textquotedblleft{}No.\textquotedblright{}}

\noindent She reached out her hand and touched the black cloth.

\noindent The cloth was warm.

\vspace{0.5em}
\noindent\textit{\textbf{↪ Routes:} \rota{LUC-03} \rotasep \rota{DOC-01} \rotasep \rota{DOC-20}}
